\documentclass[11pt,a4paper]{article}
\usepackage[T1]{fontenc}
\usepackage{lmodern,amsmath,amssymb,amsthm,microtype}
\usepackage[margin=27mm]{geometry}
\usepackage[colorlinks=true,linkcolor=blue,citecolor=blue,urlcolor=blue]{hyperref}
\newtheorem{proposition}{Proposition}
\newtheorem{theorem}[proposition]{Theorem}
\newcommand{\Q}{\mathbb Q}
\setlength{\emergencystretch}{2em}
\title{A coupled one-coordinate slice of the intrinsic pole surface}
\author{}\date{17 November 2026}
\begin{document}\maketitle

For $d\in\Q$, the intrinsic pole equation is
\begin{equation}\label{eq:intrinsic}
 Y^2=d(d+v^2)^2s^4-d(d-v^2)^2(1-d)^2.
\end{equation}
Its complete prepared-input construction retains the two parameters
$v,s$. The following coupling retains only one parameter but is not
complete on all prepared inputs.

\begin{proposition}\label{prop:slice}
Let $d>0$ and put $e=(d-v^2)/(d+v^2)$. The slice
\begin{equation}\label{eq:coupling}
 s^2=(1+d)e
\end{equation}
has the explicit ordinate $Y=2d(d-v^2)$ in \eqref{eq:intrinsic}.
The rational points of this slice are in bijection with the rational
points of
\begin{equation}\label{eq:plus}
 y^2=d\bigl((1+d)^2-s^4\bigr),
\end{equation}
by
\[
 y=(1+d+s^2)v,\qquad v=\frac{y}{1+d+s^2}.
\]
The equation \eqref{eq:plus} has total degree five and coefficient
height two. Every rational solution excludes negative odd valuations
of $d$ at every prime, including two.
\end{proposition}
\begin{proof}
Under \eqref{eq:coupling}, the right side of \eqref{eq:intrinsic} is
\[
 d(d+v^2)^2 e^2\bigl((1+d)^2-(1-d)^2\bigr)
 =4d^2(d-v^2)^2.
\]
Rearranging the coupling gives
\[
 v^2=\frac{d(1+d-s^2)}{1+d+s^2}.
\]
The asserted change of variables follows; its denominator is positive.
Now suppose $\nu=v_p(d)<0$ is odd. Then $v_p(1+d)=\nu$.
The two terms on the right side of \eqref{eq:plus} have valuations
$3\nu$ and $\nu+4v_p(s)$. If both are finite, they are odd and
distinct: equality would imply $\nu=2v_p(s)$. If $s=0$, only the
first term remains. Their sum is consequently nonzero of odd
valuation and cannot be a rational square.
\end{proof}

Call $d$ prepared if
\[
 0<d<1,\qquad \operatorname{den}(d)\text{ is a square},\qquad
 d\in7\Q_2^{\times2},\qquad v_2(d)\ge4.
\]
\begin{theorem}\label{thm:counterexample}
The number
\[
 d_0=\frac{448}{13225}=7\left(\frac8{115}\right)^2
\]
is prepared, but \eqref{eq:plus} has no rational point for $d=d_0$.
Its smooth projective model has a point over every completion of
$\Q$.
\end{theorem}
\begin{proof}
The displayed fraction is reduced, $13225=115^2$, and
$v_2(d_0)=6$. Also
\[
 1+d_0=113\left(\frac{11}{115}\right)^2.
\]
The invertible rational changes
\[
 s=\frac{11}{115}t,\qquad
 y=\frac{8\cdot11^2}{115^3}z
\]
identify \eqref{eq:plus} with
\begin{equation}\label{eq:sha}
 z^2=7(113^2-t^4).
\end{equation}
The complete descent certificate in \cite{Sha} proves that this
curve is everywhere locally soluble and has no rational point.
Explicitly, a rational point with $tz\ne0$ would give
\[
 (X,Y)=(-7t^2,7tz)\in E_{791}(\Q),\qquad
 E_{791}:Y^2=X^3-791^2X,
\]
with abscissa square class $[-7]$. That certificate proves
\[
 \operatorname{rank}E_{791}(\Q)=1,
 \qquad
 \delta(E_{791}(\Q))=\{[1],[-1],[791],[-791]\}.
\]
It uses the three-dimensional pure $2$-Selmer group, an explicitly
nonzero Cassels pairing, and the point $(-784,-2940)$; it does not
infer the rank from a search. The omitted affine cases $t=0$ or
$z=0$ are impossible since $7$ and $113$ are nonsquares. The leading
coefficient $-7$ excludes rational points at infinity.
\end{proof}

\paragraph{Scope.}
This is a counterexample to uniform completeness of the specific
coupling \eqref{eq:coupling} on all prepared inputs. It is not asserted
that $d_0$ occurs in a selected actual detector tuple. The result does
not exclude a different coupled relation, a finite menu containing
other relations, or a construction tailored to actual detector
values. The original two-parameter pole equation remains complete
on the prepared class by its simultaneous-twist theorem.

\begin{thebibliography}{9}
\bibitem{Sha}
\raggedright
\emph{A prepared input with a nontrivial one-coordinate covering
obstruction}, project manuscript, 17 November 2026,
\texttt{negative\_results/papers/pole\_parameters/}\allowbreak
\texttt{prepared\_sha\_obstruction\_20261117\_v57.tex};
exact checkers and Cassels certificate in
\texttt{negative\_results/computation/prepared\_sha\_v57/}.
Its primary descent references are Z.~Wang,
\emph{Congruent elliptic curves with non-trivial Shafarevich--Tate
groups}, arXiv:1511.03810, Sections 2.1--2.2, and S.~Zhang,
\emph{On non-congruent numbers as multiples of non-congruent numbers},
28 August 2025, Sections 2.1--2.2.
\end{thebibliography}
\end{document}
